% Independent continuation; the main parity claim is not promoted to a theorem.
\subsection*{The remaining cubic comparison after the quadratic pentagon equation}

\paragraph{MB.1 (Proved): an entirely finite formula for the residual duplication character.}
Keep the intrinsic Schreier convention of MD.1. Put
$p=(123)$, $t=(1234)$ in $S_4$, and
$P=(0,\omega)$, $Y=(1,1)$ in
$\mathbb F_4\rtimes C_3\simeq A_4$.
For every actual $\Phi=(1,f)\in\mathcal H$, write
\[
 f(p,t)=t^{2\nu_4(\Phi)}p^{a_f(\Phi)},\qquad
 f(P,Y)=(k_f,1).
\]
FP.37 proves the first normal form and its additive coordinate law.
KM.3 gives $k_f\in\mathbb F_2$ and
$k_f=\Psi_3(\Phi)/2\pmod2$.
Consequently
\[
 (a_f,\nu_4,k_f):\mathcal H\longrightarrow
           C_3\times C_2\times C_2
\]
is a continuous homomorphism, and
\[
 \xi_4:=\nu_4+k_f
\]
is a literal finite-evaluation character with no gamma-function
normalization in its definition. The source word evaluations lie
in $S_4\times A_4$; MD.9 proves that their joint word map is onto
that product. No surjectivity of the displayed gentle-group
homomorphism is asserted.

For genuine $\Phi\in\mathcal H_5$, FP.37 and MD.4 imply
\[
 d_2(\Phi)=\xi_4(\Phi),\qquad
 \Delta(\Phi)=\kappa_d(\Phi)+\nu_4(\Phi)\pmod2.
 \tag{MB.1}
\]
Thus the exact outstanding assertion is
\[
 \boxed{\quad
 f_\Phi(p,t)\in t^{2\kappa_d(\Phi)}\langle p\rangle
 \qquad(\Phi\in\mathcal H_5).
 \quad}
 \tag{MB.2}
\]
Equivalently, the two continuous characters $\kappa_d\bmod2$
and $\nu_4$ agree. Proving $\xi_4=0$ alone is insufficient:
MD.10 supplies an actual gentle element with the entire
beta/Kummer coefficient package trivial but nontrivial cubic defect.

On either fixed locus $C_L\cap\mathcal H_5$ or
$C_P\cap\mathcal H_5$, MD.8 and FP.37 already give
$\nu_4=\Psi_3/2\pmod2$. There the unresolved assertion is
exactly $\kappa_d=\Psi_3/2\pmod2$.

\paragraph{MB.2 (Known source boundary): the Mordell comparison is still arithmetic.}
In Nakamura--Tsunogai--Yasuda (2010), Proposition 4.4 defines
the triangle-group action on the Mordell Tate module for every
full profinite GT element. The comparison with the standard
braid action is arithmetic: Propositions 4.11 and 4.13 and
Theorem 4.16 quantify over $G_{\mathbb Q}$. The introduction,
printed p.432, expressly leaves extension of Proposition 4.13
to all GT unknown. Thus these results do not prove (MB.2).
The missing input is compatibility of the two specified actions
through the Mordell-locus path, with its signed normalization;
the existence of the two actions separately does not supply it.

\paragraph{MB.3 (An action detector, with a precise pentagon limitation).}
Any next comparison must target the one scalar ambiguity
in the cubic frame, instead of another beta factorization identity.
FP.37 identifies an actual source element with Burau value one,
abelianization $e_5$ modulo two, and forgotten signed residue
$[5I]$. Its image centralizes the cubic torsion in that finite
target, so that target cannot recover $\kappa_d\bmod2$.

RH.2 provides a more specific replacement for the earlier enormous
characteristic core: the relation module of the level-twelve
cover, reduced modulo $576$, together with the finite extension
$\Gamma_{576}$. Its cohomology calculation kills the relevant
automorphism cocycle after reduction to $\Gamma_6$, and FP.23's
permutation-module extension removes the centralizer projection
ambiguity. Induce that specified representation from the
label-five stabilizer to $G_5$, retaining the actual fifth-point
action and the literal $S_4$ evaluation defining $\nu_4$.

This gives an actual GT-stable action detector: if $L$ is the
label-five stabilizer, then $q_5^{-1}(N'N^6)$ is stable under
the specified genuine GT action on $L$. Its normal core in
$G_5$ is therefore GT-stable too. The resulting finite quotient
retains the forgotten $\Gamma_6$ action, whose kernel on
$\mathcal H$ is contained in $\ker\Delta$ by FP.23's two
centralizer-projection arguments.

However FP.30 supplies an essential obstruction to this proposed
next calculation: a representation induced solely from forgetting
points kills the two-cycle pentagon defect. Thus this detector
cannot itself furnish the missing parity equation merely by
evaluating the pentagon in it. A successful relation-module or
Fox-derivative attack must additionally retain a genuine
five-point component not factoring through those forgetful maps.
No such compatibility calculation is claimed here. The target
remains to derive (MB.2) for actual GT actions, and a finite
assignment alone would not supply a global lift.
